\chapter{Diskussion}
\label{cha:Diskussion}

Dieses Kapitel interpretiert die in Kapitel~\ref{cha:Ergebnisse} berichteten
Ergebnisse im Kontext der vier Nebenfragestellungen. Jeder Abschnitt schließt
mit einer expliziten Antwort auf die zugehörige Fragestellung.

\section{Leistungsfähigkeit in der Trainingsumgebung}
\label{sec:disk_nf1}

\subsection{Trainingsverhalten und Konvergenz}

Die zwei beobachteten Phasen im Belohnungsverlauf lassen sich auf die
Belohnungsstruktur zurückführen, die dichte Komponenten mit
ereignisbasierten Boni kombiniert. In der Konvergenzphase (Iteration~0
bis~441) ermöglichen die dichten Komponenten ($r_{\mathrm{prog}}$,
$r_{\mathrm{nah}}$, $r_{\mathrm{prox}}$) den initialen Anstieg, indem
sie bereits für Teilfortschritte Signal liefern. Ab circa Iteration~200
setzt zusätzlich die Erfolgsbelohnung $r_{\mathrm{erf}}$ ein, sobald der
Agent erstmals konsistente Landungen erreicht, und beschleunigt den
Belohnungsanstieg. Ohne die dichten Komponenten würde dem Agenten in der
frühen Trainingsphase das Gradientensignal für Teilfortschritte fehlen,
was die Konvergenz vermutlich erheblich verzögern würde. Der abrupte Kollaps ab Iteration~442 betrifft ausschließlich die
Mittelwerte der Strategie bei nahezu unveränderter Standardabweichung
($0{,}193 \rightarrow 0{,}194$), was auf eine katastrophale
Strategieaktualisierung hindeutet, nicht auf einen graduellen
Stabilitätsverlust. Obwohl der Clipping-Mechanismus von
PPO~\autocite{schulmanProximalPolicyOptimization2017} große
Strategieänderungen pro Iteration begrenzen soll, reicht er nicht aus,
um solche abrupten Einbrüche zu verhindern. Eine mögliche Erklärung
bieten \textcite{moallaNoRepresentationNoTrust2024}, die zeigen, dass
eine schleichende Verschlechterung des Feature-Ranks im
Actor-Netzwerk den Clipping-Mechanismus unterlaufen und zu
unumkehrbaren Leistungseinbrüchen führen kann, wobei die
Standardabweichung als globaler Parameter unbeeinflusst bleibt. Ob
dieser Mechanismus den beobachteten Kollaps erklärt, lässt sich ohne
Analyse des Feature-Ranks nicht abschließend beurteilen. % REVIEW: Moalla et al. 2024 (NeurIPS) — Quelle lesen und bestätigen

Die im Trainingsverlauf ansteigende Landezeit lässt sich auf die
Struktur der Belohnungsfunktion zurückführen. Die beiden dominanten
Komponenten Fortschritt $r_{\mathrm{prog}}$ (Gewicht~$1{,}0$) und Nähe
$r_{\mathrm{nah}}$ (Gewicht~$2{,}0$) werden bei steilem Sinkflug
maximiert, da die 3D-Distanz zum Ziel pro Schritt am stärksten abnimmt.
Entscheidend ist dabei die unterschiedliche Natur beider Komponenten:
$r_{\mathrm{prog}} = d_{t} - d_{t+1}$ misst die Distanzverringerung pro
Schritt und ist damit potentialbasiert im Sinne von
\textcite{ngPolicyInvarianceReward1999}; solches Reward Shaping
verändert die optimale Strategie nicht.
$r_{\mathrm{nah}} = \exp(-d_{t+1})$ hingegen ist nicht potentialbasiert,
sondern ein zustandsabhängiger Pro-Schritt-Reward: Jeder Zeitschritt in
Zielnähe liefert Belohnung, unabhängig davon, ob der Agent sich dem Ziel
weiter nähert. Da die Zeitstrafe deaktiviert ist (Gewicht~$0{,}0$),
existiert kein Gegengewicht zu diesem Verweil-Anreiz. Dieses Phänomen
ist in der Literatur bekannt: \textcite{randlovLearningDriveBicycle1998}
dokumentieren, dass ein nicht-potentialbasierter Shaping-Reward einen
Agenten dazu bringt, Kreise zu fahren statt das Ziel zu erreichen, da
die akkumulierte Zwischenbelohnung die Zielbelohnung übersteigt. % REVIEW: Randlov & Alstrom 1998 (ICML) — Quelle lesen und bestätigen
Der Agent lernt folglich eine Strategie, die steilen Abstieg mit
ausgedehntem Schweben in Zielnähe kombiniert, was die beobachtete
Zunahme der Episodenlänge im Trainingsverlauf erklärt.

\subsection{Dominante Fehlerart und Flugverhalten}

Die Failure-Mode-Analyse zeigt ein eindeutiges Bild:
Hinderniskollisionen verursachen 28,8\,\% aller Episodenausgänge und
stellen damit die einzig relevante Fehlerart dar. Zeitüberschreitungen und
Bodenkollisionen treten nicht auf (jeweils 0,0\,\%); der Agent erreicht
in jeder Episode die Zielumgebung. Die Nachuntersuchung mit
oberflächenbasierter Kollisionserkennung
(Abschnitt~\ref{subsec:phase1_mesh_contacts}) relativiert dieses
Ergebnis: Mit geometrisch exakter Kontaktprüfung steigt die Erfolgsrate
auf 98,0\,\%, was zeigt, dass der Großteil der registrierten Kollisionen
auf die konservative AABB-Metrik zurückgeht und nicht auf tatsächliche
Berührungen der Hindernisoberfläche. Der limitierende Faktor der
berichteten 71,0\,\% Erfolgsrate ist somit primär die Kollisionsmetrik,
nicht das erlernte Flugverhalten.

Zur Einordnung: \textcite{bartolomeiAutonomousEmergencyLanding2022}
berichten vergleichbare Erfolgsraten (über 70\,\%) für ein RL-System mit
Tiefenbildeingabe und diskretem Aktionsraum. Deren Aufgabe ist jedoch
die Selektion sicherer Landezonen in offenem Gelände, nicht die
Navigation durch Hindernisse zu einem festen Ziel. Dass der hier
trainierte Agent mit kontinuierlichem Aktionsraum in einer
Korridorumgebung mit erzwungenen Ausweichmanövern vergleichbare Raten
erreicht, spricht für die Leistungsfähigkeit der gewählten Architektur.

Das qualitative Flugverhalten ist reaktiv: Der Agent nähert sich einem
Hindernis auf direktem Weg, bremst auf kurzer Distanz ab, weicht lateral
aus und setzt den Abstieg fort. Dieses Muster ist eine direkte Konsequenz
des begrenzten lokalen Sensorhorizonts: Die nach unten gerichtete Kamera
erfasst Hindernisse erst bei relativer Nähe, sodass vorausschauende
Planung nicht möglich ist. \textcite{chenDRQN3DObstacleAvoidance2021}
zeigen, dass Agenten mit eingeschränktem Sichtfeld ohne temporales
Gedächtnis zwangsläufig reaktives Ausweichverhalten im letzten Moment
entwickeln; erst eine LSTM-Erweiterung ermöglicht vorausschauendes
Handeln. Die hier eingesetzte CNN-MLP-Architektur ohne rekurrente
Komponente ist daher strukturell auf reaktives Verhalten beschränkt. % REVIEW: Chen et al. 2021 (IROS) — Quelle lesen und bestätigen Bemerkenswert ist die
Präzision, mit der Hindernisse in geringem Abstand passiert werden
(mittlerer Minimalabstand 0,52\,m bei einem Kollisionsradius von
0,3\,m).

\subsection{Perceptual Aliasing}

Ein fundamentales Problem der gewählten Sensorarchitektur zeigt sich
unmittelbar in den aufgezeichneten Tiefenbildern: Boden und Hindernisse
erzeugen bei Annäherung nahezu identische Tiefensignaturen. Während die
Annäherung an den Boden das Ziel darstellt, muss jene an Hindernisse
vermieden werden. Der Agent löst diese Ambiguität implizit: Die relative
Positionsinformation im Zustandsvektor kodiert die Entfernung zum Ziel und
damit indirekt die erwartete Flughöhe, sodass das Modell lernen kann, ob
eine abnehmende Tiefe Boden oder Hindernis bedeutet.
\textcite{bartolomeiAutonomousEmergencyLanding2022} adressieren dieses
Problem, indem sie neben Tiefenbildern zusätzlich semantische
Segmentierungen als Eingabe verwenden, die eine explizite Unterscheidung
von Oberflächentypen ermöglichen. Im vorliegenden System steht eine
solche Segmentierung nicht zur Verfügung; die Disambiguierung stützt
sich ausschließlich auf den Zustandsvektor.

\subsection{Beantwortung NF1}

Der trainierte Agent erreicht in der Korridorumgebung eine Erfolgsrate
von 71,0\,\% unter der konservativen AABB-Kollisionsmetrik; mit
oberflächenbasierter Kontakterkennung steigt dieser Wert auf 98,0\,\%.
Hinderniskollisionen (28,8\,\%) sind die einzig relevante Fehlerart,
wobei der Großteil auf die Konservativität der AABB-Metrik
zurückzuführen ist. Der Agent erlernt eine reaktive Ausweichstrategie,
die Hindernisse mit hoher Präzision passiert. Ihre Leistungsgrenze wird
primär durch den begrenzten Sensorhorizont der nach unten gerichteten
Tiefenkamera bestimmt. Die nicht-potentialbasierte
Nähebelohnung $r_{\mathrm{nah}}$ erzeugt dabei einen Verweil-Anreiz in
Zielnähe, der die Landezeiten erhöht, aber die geforderte Landepräzision
begünstigt.

\section{Transferleistung auf den Weinberg}
\label{sec:disk_nf2}

\subsection{Success-Drop und Failure-Mode-Shift}

Der Leistungsabfall von Phase~1 zu Phase~2 ist erwartbar: Der Agent wurde
ausschließlich in der Korridorumgebung trainiert und sieht die
Weinberggeometrie erstmals zur Evaluationszeit. Aufschlussreich ist weniger
der Drop selbst als dessen Zusammensetzung.
% TODO: Interpretation je nach dominantem Failure-Mode einfügen.
% Falls Timeouts dominieren: Ein erhöhter Timeout-Anteil deutet darauf hin,
% dass der Agent sich dem Ziel zwar nähert, aber in der Nähe der Landefläche
% kein stabiles Hovern erreicht.
% Falls Terrain-Kollisionen dominieren: Die Verschiebung von
% Hinderniskollisionen zu Terrain-Kollisionen spiegelt den geometrischen
% Unterschied der Umgebungen wider.

\subsection{Einordnung als Sim-to-Sim Transfer}

Die Phase-2-Evaluation ist als Zero-Shot Sim-to-Sim Transfer einzuordnen:
Der Agent wird ohne Fine-Tuning und ohne Domain Randomization in einer
geometrisch verschiedenartigen Umgebung evaluiert. Dass überhaupt eine
partielle Übertragbarkeit besteht, lässt sich auf die umgebungsagnostische
Beobachtungsstruktur zurückführen. Der Zustandsvektor kodiert relative
Positionen, und das Tiefenbild enthält keine texturbasierten Merkmale, die
an die Korridorgeometrie gebunden wären. Diese Einordnung sollte nicht
überbewertet werden: Sim-to-Sim Transfer unter kontrollierten Bedingungen
ist ein schwächeres Ergebnis als Sim-to-Real Transfer oder Generalisierung
über diverse Trainingsumgebungen.

\subsection{Perceptual Aliasing im Weinberg}

Das in Abschnitt~\ref{sec:disk_nf1} beschriebene Perceptual-Aliasing-Problem
bietet eine Erklärung für die erhöhte Terrain-Kollisionsrate im Weinberg. Im
Korridor ist die Landefläche hindernisfrei: Sobald der Agent die letzte
Hindernisschicht passiert hat, kann er das Tiefenbild in der terminalen Phase
gefahrlos ignorieren und sich ausschließlich am Zustandsvektor orientieren.
Dass der Agent genau diese Strategie erlernt, ist plausibel, da das
Tiefenbild in Bodennähe ohnehin nur noch abnehmendes Bodenrelief zeigt und
keinen handlungsrelevanten Informationsgewinn bietet. Im Weinberg bricht
diese Strategie: Rebenreihen ragen als vertikale Strukturen direkt in die
Landezone, sodass die Tiefenwahrnehmung gerade in der terminalen Phase
entscheidungsrelevant bleibt. Ein Agent, der gelernt hat, das Tiefenbild in
Zielnähe zu ignorieren, kollidiert mit Vegetation, die er bei aktiver
Auswertung hätte umgehen können.

Ein möglicher Lösungsansatz wäre die Ergänzung der Beobachtung um eine
semantische Segmentierung, die Boden, Hindernisse und Vegetation als separate
Klassen ausweist. Anstatt die Boden-Hindernis-Unterscheidung implizit aus dem
Tiefenbild ableiten zu müssen, könnte der Agent direkt lernen, welche
Strukturen gefahrlos überflogen werden dürfen und welche ein Ausweichmanöver
erfordern.

\subsection{Beantwortung NF2}

% TODO: Explizite Antwort auf NF2 einfügen. Beispielstruktur:
% Die erlernte Navigationsleistung lässt sich partiell auf die
% Weinbergumgebung übertragen: Die Erfolgsrate sinkt von [X\,\%] auf
% [Y\,\%], wobei sich der dominante Failure-Mode von [A] zu [B]
% verschiebt. Der Transfer gelingt ohne Fine-Tuning, ist aber durch das
% verschärfte Perceptual-Aliasing-Problem in der terminalen Landephase
% begrenzt.

\section{Leistungslücke zum Pfadplaner}
\label{sec:disk_nf3}

\subsection{Inferenzzeit und Echtzeitfähigkeit}

Ein entscheidender praktischer Vorteil des RL-Agenten liegt in der
Inferenzzeit. Ein Forward-Pass durch das Netzwerk benötigt konstante Zeit
unabhängig von der Umgebungskomplexität, während RRT-Connect iterativ Pfade
sampelt und Kollisionsprüfungen auswertet. Auf eingebetteter Hardware, wie
sie für reale Drohnenanwendungen relevant wäre, ist der RL-Agent
echtzeitfähig, während RRT-Connect entweder ein stark begrenztes Sampling-Budget
oder leistungsfähigere Hardware erfordert.

\subsection{Landegeschwindigkeit}

Die in Abschnitt~\ref{subsec:meth_versuch1} beschriebene
Informationsasymmetrie lässt eine deutliche Leistungslücke zugunsten des
RRT-Connect-Planers erwarten. Dass der RL-Agent dennoch vergleichbare Erfolgsraten
erzielt, zeigt, dass die implizite Hindernisrepräsentation im Tiefenbild
für diese Aufgabe ausreichend ist. Der RL-Agent erzielt zudem kürzere
Landezeiten: Er lernt direkte Trajektorien, die keine konservativen
Sicherheitsabstände einhalten, wie sie ein sampling-basierter Planer
typischerweise produziert.

\subsection{Planungshorizont}

RRT-Connect plant global: Der gesamte Pfad wird vor Flugbeginn berechnet. Der
RL-Agent entscheidet lokal, pro Zeitschritt, auf Basis der aktuellen
Beobachtung. In dieser Arbeit sind alle Hindernisse statisch, sodass der
globale Planungshorizont von RRT-Connect ein reiner Vorteil ist. In dynamischen
Umgebungen kehrt sich dieses Verhältnis um: RRT-Connect müsste bei jeder Änderung
replanen, während der RL-Agent durch seine lokale Entscheidungsstruktur
inhärent auf veränderte Bedingungen reagieren kann. Dieses Argument bleibt
im Kontext dieser Arbeit theoretisch, ist aber für die praktische
Motivation des lernbasierten Ansatzes relevant.

\subsection{Beantwortung NF3}

% TODO: Explizite Antwort auf NF3 einfügen. Beispielstruktur:
% Die Leistungslücke zwischen RL-Agent und RRT-Connect-Pfadplaner beträgt
% [X Prozentpunkte] in der Erfolgsrate. Der RL-Agent kompensiert diese
% Differenz durch kürzere Landezeiten und Echtzeitfähigkeit.

\section{Beitrag der Beobachtungskomponenten}
\label{sec:disk_nf4}

Eine quantitative Beantwortung von NF4 erfordert eine Ablationsstudie,
in der ein State-only-Modell (ohne Tiefenbild) unter identischen
Bedingungen trainiert und evaluiert wird. Eine solche Ablation wurde im
Rahmen dieser Arbeit nicht durchgeführt; der individuelle Beitrag von
Tiefenwahrnehmung und propriozeptivem Zustandsvektor lässt sich daher
nur indirekt abschätzen.

Das beobachtete Flugverhalten liefert qualitative Hinweise: Die
lateralen Ausweichmanöver vor Hindernissen
(Abschnitt~\ref{sec:phase1}) setzen Hindernisinformation voraus, die
ausschließlich im Tiefenbild enthalten ist, da der Zustandsvektor keine
Hindernispositionen kodiert. Ohne visuelle Eingabe könnte der Agent
lediglich den direkten Abstieg zum Ziel erlernen; erfolgreiche Episoden
wären auf Hinderniskonfigurationen beschränkt, die zufällig einen freien
Korridor auf der Abstiegslinie belassen. Die Transferleistung auf
unbekannte Hindernisformen (Phase~1) und auf die Weinbergumgebung
(Phase~2) stützt die Hypothese, dass das CNN geometrische Strukturen
erkennt, die über die Trainingsumgebung hinaus generalisieren. Welche
konkreten Merkmale extrahiert werden, bleibt ohne
Feature-Visualisierung (z.\,B.\ Grad-CAM) offen.

\subsection{Beantwortung NF4}

Die Frage nach dem individuellen Beitrag der Beobachtungskomponenten
kann ohne Ablationsstudie nicht quantitativ beantwortet werden. Die
qualitative Analyse des Flugverhaltens legt nahe, dass die
Tiefenwahrnehmung die Hindernisumgehung ermöglicht, während der
Zustandsvektor die Zielnavigation und den kontrollierten Abstieg
steuert. Eine experimentelle Validierung dieser Hypothese bleibt
zukünftiger Arbeit vorbehalten
(vgl.\ Abschnitt~\ref{sec:disk_limitationen}).

\section{Limitationen}
\label{sec:disk_limitationen}

\textbf{Einzelner Seed.}
Alle Ergebnisse basieren auf einem einzigen Trainingslauf. Die
Bootstrap-Konfidenzintervalle quantifizieren die Evaluationsvarianz über
Hinderniskonfigurationen und Startpositionen, nicht die Trainingsvarianz.
Ein anderer Seed könnte zu qualitativ anderem Verhalten führen. Die
berichteten Erfolgsraten sind daher als Punktschätzung eines einzelnen
Trainingslaufs zu verstehen, nicht als robuste Aussage über die Methode.

\textbf{Sim-to-Real Gap.}
Die Simulation bietet perfekte Sensorik ohne Rauschen, keine
Windeinflüsse, keine Motorverzögerung und keine Kommunikationslatenz.
Ergebnisse sind nicht direkt auf reale Drohnen übertragbar. Es wurde kein
Domain-Randomization-Ansatz implementiert, der die Robustheit gegenüber
sensorischen Störungen erhöhen könnte.

\textbf{Reward Shaping.}
Die Belohnungsgewichte wurden manuell in Vorversuchen abgestimmt. Andere
Gewichtungen könnten zu qualitativ anderem Verhalten führen. Ein
systematisches Hyperparameter-Tuning der Reward-Gewichte wurde nicht
durchgeführt.

\textbf{Feste Zielposition.}
Das Landeziel ist in allen Episoden identisch. Obwohl die Beobachtung
relativ kodiert ist, könnte der Agent auf positionsspezifische Muster
überangepasst sein.

\textbf{Statische Umgebung.}
Alle Hindernisse sind statisch. Reale Weinberge weisen bewegliche
Vegetation, Wind und andere dynamische Elemente auf.

\textbf{Keine Ablation und keine Feature-Visualisierung.}
Der individuelle Beitrag von Tiefenwahrnehmung und Zustandsvektor wurde
nicht durch eine Ablationsstudie quantifiziert
(vgl.\ Abschnitt~\ref{sec:disk_nf4}). Die Schlussfolgerungen über die
gelernten CNN-Repräsentationen sind indirekt aus Trajektorien und
Transferleistung abgeleitet. Eine Grad-CAM-Analyse oder vergleichbare
Methode würde direktere Evidenz liefern.
